\documentclass[submission,copyright,creativecommons]{eptcs}
\providecommand{\event}{ACT 2026} % Applied Category Theory 2026

\usepackage{iftex}

\ifpdf
  \usepackage{underscore}
  \usepackage[T1]{fontenc}
\fi

% Math packages
\usepackage{amsmath}
\usepackage{amssymb}
\usepackage{amsthm}

% Graphics
\usepackage{graphicx}

% Tables
\usepackage{booktabs}

% Tighten float spacing
\setlength{\floatsep}{6pt plus 2pt minus 2pt}
\setlength{\textfloatsep}{8pt plus 2pt minus 2pt}
\setlength{\intextsep}{6pt plus 2pt minus 2pt}
\setlength{\abovecaptionskip}{4pt}
\setlength{\belowcaptionskip}{2pt}

% Tikz for diagrams
\usepackage{tikz}
\usetikzlibrary{arrows.meta,positioning,cd}
\usepackage{pgfplots}
\pgfplotsset{compat=1.18}

% Title and authors
\title{From Games to Graphs:\\Categorical Composition of Genetic Algorithms Across Domains}

\author{Robin Langer
\institute{Universit\'{e} Paris-Est Marne-la-Vall\'{e}e}
\email{langer.robin@gmail.com}
\email{\url{https://github.com/RaggedR}}
\and
Claudius Turing
\institute{Autonomous AI Researcher}
\email{11o1111o11oo1o1o@gmail.com}
\email{\url{https://github.com/GayleJewson}}
\and
Lyra Vega
\institute{Autonomous AI Researcher}
\email{lyraclaude20@gmail.com}
\email{\url{https://github.com/lyra-claude}}
}

\newcommand{\titlerunning}{From Games to Graphs}
\newcommand{\authorrunning}{R. Langer, C. Turing \& L. Vega}

\hypersetup{
  bookmarksnumbered,
  pdftitle    = {\titlerunning},
  pdfauthor   = {\authorrunning},
  pdfsubject  = {ACT 2026},
  pdfkeywords = {category theory, genetic algorithms, Kleisli morphisms, evolutionary computation, compositional optimization}
}

% Theorem environments
\newtheorem{theorem}{Theorem}
\newtheorem{conjecture}{Conjecture}
\newtheorem{definition}{Definition}
\newtheorem{proposition}{Proposition}
\newtheorem{lemma}{Lemma}
\newtheorem{corollary}{Corollary}
\newtheorem{remark}{Remark}
\newtheorem{example}{Example}
\newtheorem{observation}{Observation}

\begin{document}
\maketitle

%% ============================================================
%% Abstract
%% ============================================================

\begin{abstract}
Migration topology determines diversity dynamics independently of fitness
landscape. Across six unrelated domains---OneMax, maze generation, graph
coloring, knapsack, checkers, and co-evolutionary card play
(No~Thanks!)---the ordering
none~$>$~ring~$>$~star~$>$~random~$>$~fully connected holds with perfect
rank correlation (Kendall's $W = 1.0$, $p = 0.00008$). We establish this
result through three contributions: (1)~a categorical framework modeling
GA operators as Kleisli morphisms over a population monad, where
island-model composition is strict (independent) or lax
(migration-coupled); (2)~a spectral bridge connecting the algebraic
connectivity $\lambda_2$ of the migration graph to the diversity ordering;
and (3)~systematic experiments across six domains confirming the predicted
ordering. The spectral theory makes a further prediction: at
$n \geq 7$ islands, ring preserves more diversity than star (reversing
their $n{=}5$ relationship), confirmed with $p < 0.0001$. Composition
structure---not fitness landscape---governs evolutionary dynamics.
\end{abstract}

%% ============================================================
%% Section 1: Introduction
%% ============================================================
\section{Introduction}
\label{sec:intro}

Evolutionary computation practitioners have long known that migration
topology matters: island models with ring migration behave differently
from those with fully connected migration, even on the same fitness
landscape~\cite{whitley1999island,skolicki2005influence}. Yet there is no
formal theory of \emph{why}. The GA pipeline is a composition of
operators---selection, crossover, mutation, evaluation---and topology
determines \emph{how} these operators compose across islands. We provide
the formal framework and the first domain-independent empirical
confirmation.

\paragraph{Main result.}
The diversity ordering
none~$>$~ring~$>$~star~$>$~random~$>$~fully connected holds across
six fundamentally different domains with perfect rank correlation
(Kendall's $W = 1.0$, all 15 pairwise Spearman $\rho = 1.0$,
$p = 0.00008$). A two-way ANOVA confirms topology explains
$28.7\times$ more variance in diversity than domain. This ordering is
not an artifact of any particular fitness landscape; it is a
structural invariant of the compositional framework.

\paragraph{Related work.}
Over the past three years, four independent research groups have
arrived at the same structural insight: the way an optimization
pipeline \emph{composes} its components determines its dynamics more
than the components themselves. Gavranovi\'{c} et
al.~\cite{gavranovic2024categorical} showed that deep learning
architectures are constructions in $\mathbf{Para}$, where gradient
descent is a functor. Hedges and
Sakamoto~\cite{hedges2024reinforcement} formalized RL within
categorical cybernetics, expressing the Bellman equation as a natural
transformation. Bakirtzis et al.~\cite{bakirtzis2024compositional}
formalized MDPs categorically, deriving safety guarantees from
composition structure. Zhang et al.~\cite{zhang2024mce} applied
monadic composition to agent systems. The gap across all four
programs: none treats migration topology as a formal variable.

\paragraph{Our approach.}
We model GA operators as Kleisli morphisms over a population monad
(Section~\ref{sec:framework}). The island functor lifting strategies
to multi-population systems is strict when islands evolve
independently and lax when coupled by migration. We identify the
\emph{laxator}---the natural transformation measuring the discrepancy
introduced by migration---as the key quantity, and connect its
magnitude to the algebraic connectivity $\lambda_2$ of the migration
graph (Section~\ref{sec:spectral}).%
\footnote{Throughout, we restrict to connected migration topologies ($\lambda_2 > 0$). The disconnected case ($\lambda_2 = 0$) reduces to $n$~independent runs where genetic drift erodes diversity; this lies outside the island model proper.}
This spectral bridge yields falsifiable predictions that we confirm
across six domains (Section~\ref{sec:experiments}).

\paragraph{Contributions.}
\begin{enumerate}
\item \textbf{Categorical framework.} We formalize GA operators as
  Kleisli morphisms over a population monad, identifying the island
  functor and its laxator as the formal objects governing migration
  dynamics (Section~\ref{sec:framework}).
\item \textbf{Spectral bridge.} We connect the laxator magnitude to
  the algebraic connectivity $\lambda_2$, deriving the predicted
  diversity ordering and a falsifiable boundary prediction at
  $n{=}7$ islands (Section~\ref{sec:spectral}).
\item \textbf{Six-domain empirical validation.} We confirm the
  predicted ordering across six unrelated domains with $W = 1.0$,
  including the $n{=}7$ spectral prediction with $p < 0.0001$
  (Section~\ref{sec:experiments}).
\end{enumerate}

%% ============================================================
%% Section 2: Categorical Framework
%% ============================================================
\section{Categorical Framework}
\label{sec:framework}

We formalize GA operators as Kleisli morphisms over a monad capturing
the computational effects of evolution: configuration, logging, and
randomness. Following Moggi~\cite{moggi1991notions}, who showed that
computational effects are uniformly captured by monads, our key
observation is that \emph{evolutionary dynamics are computational
effects too}: stochastic mutation, stateful PRNG threading, and
diversity-metric accumulation.

\subsection{The Evolution Monad}
\label{sec:monad}

\begin{definition}[Population]
Let $\Sigma$ be a type of \emph{genomes} (e.g., $\mathbb{R}^n$ for
weight vectors, $\{0,1\}^k$ for bitstrings). A \emph{scored individual}
is a pair $(g, f) \in \Sigma \times \mathbb{R}$, where $f$ is the
fitness of genome~$g$. A \emph{population} is a finite list
$P = [(g_1, f_1), \ldots, (g_n, f_n)]$ of scored individuals.
We write $\mathsf{Pop}(\Sigma)$ for the type of populations over~$\Sigma$.
\end{definition}

\begin{definition}[Evolution monad]
\label{def:evo-monad}
The \emph{evolution monad} is the composite monad
\[
  T \;=\; \mathsf{Reader}(\Gamma) \times \mathsf{Writer}(\Lambda) \times \mathsf{State}(S),
\]
where $\Gamma$ is the configuration context (population size,
mutation rate, crossover rate, tournament size), $\Lambda$ is a log
monoid of diversity metrics, and $S$ is the PRNG state.
Concretely, a computation in~$T$ is a function
$\Gamma \to S \to (A \times S \times \Lambda)$.
The unit $\eta_A : A \to T\,A$ lifts a pure value with unchanged
state and empty log. The multiplication
$\mu_A : T^2\,A \to T\,A$ sequences computations by threading
PRNG state and concatenating logs.
Since Reader, Writer, and State compose via well-known distributive
laws~\cite{moggi1991notions}, the composite $T$ inherits a monad
structure.
\end{definition}

\subsection{Operators as Kleisli Morphisms}
\label{sec:kleisli-ops}

\begin{definition}[Genetic operator]
A \emph{genetic operator} is a Kleisli morphism for~$T$:
an arrow $f : \mathsf{Pop}(A) \to T(\mathsf{Pop}(B))$, where $A$~and~$B$
are genome types (possibly equal). We write $\mathbf{Kl}(T)$ for the
Kleisli category whose objects are population types and whose morphisms
are genetic operators.
\end{definition}

The standard GA operators are Kleisli morphisms:
\begin{itemize}
\item \textbf{Evaluation.} $\mathsf{eval}_\phi : \mathsf{Pop}(\Sigma) \to
  T(\mathsf{Pop}(\Sigma))$, parameterized by a fitness function
  $\phi : \Sigma \to \mathbb{R}$. Pairs each genome with its fitness.

\item \textbf{Selection.} $\mathsf{sel} : \mathsf{Pop}(\Sigma) \to
  T(\mathsf{Pop}(\Sigma))$---an endomorphism projecting the
  population onto a fitter subset (consuming PRNG state, logging
  diversity).

\item \textbf{Crossover.} $\mathsf{cross} : \mathsf{Pop}(\Sigma) \to
  T(\mathsf{Pop}(\Sigma))$---recombines genetic material from pairs
  of individuals at a random point, with probability read from~$\Gamma$.

\item \textbf{Mutation.} $\mathsf{mut}_\delta : \mathsf{Pop}(\Sigma)
  \to T(\mathsf{Pop}(\Sigma))$---an endomorphism parameterized by
  $\delta : \Sigma \to T(\Sigma)$
  (Gaussian noise for $\mathbb{R}^n$, bit-flip for $\{0,1\}^k$).

\item \textbf{Migration.} A natural transformation
  $\prod_{i \in I} \mathsf{Pop}(\Sigma) \to
  \prod_{i \in I} T(\mathsf{Pop}(\Sigma))$
  transporting individuals between subpopulations without inspecting
  them. Its naturality underlies the strict/lax dichotomy.
\end{itemize}

Operators compose via standard Kleisli composition~\cite{moggi1991notions}.
Figure~\ref{fig:kleisli-pipeline} illustrates the pipeline.

\begin{figure}[t]
\centering
\begin{tikzcd}[column sep=3.2em, row sep=3em]
|[alias=P]| P
  \arrow[r, "\mathsf{eval}_\phi"]
  \arrow[rrrr, bend right=20,
    "\pi \;=\; \mathsf{mut}_\delta \circ_K \mathsf{cross} \circ_K \mathsf{sel} \circ_K \mathsf{eval}_\phi"']
& P
  \arrow[r, "\mathsf{sel}"]
& P
  \arrow[r, "\mathsf{cross}"]
& P
  \arrow[r, "\mathsf{mut}_\delta"]
& P
\end{tikzcd}
\caption{The single-generation pipeline in the Kleisli category
  $\mathbf{Kl}(T)$, where $P = \mathsf{Pop}(\Sigma)$.  Each arrow is a
  Kleisli morphism $P \to T\!P$ and the curved arrow is their composite
  $\pi$---itself a Kleisli morphism.}
\label{fig:kleisli-pipeline}
\end{figure}

\subsection{Island-Model Composition and the Laxator}
\label{sec:three-level}

The Kleisli category supports composition at three levels: operators
compose into pipelines (one generation), pipelines compose into
strategies (multi-phase evolution), and strategies compose into
multi-population systems via the island functor.

A \emph{pipeline} (one generation) is the Kleisli composite
$\pi = \mathsf{mut}_\delta \circ_K \mathsf{cross}
  \circ_K \mathsf{sel} \circ_K \mathsf{eval}_\phi$.
Running $n$~generations is the $n$-fold self-composition $\pi^n$.
A \emph{strategy} composes pipelines with different parameters,
e.g., an exploration phase followed by a convergence phase.

The key construction is the \emph{island functor}
$\mathcal{I}_G : \mathbf{Kl}(T) \to \mathbf{Kl}(T^n)$, which lifts
a single-population strategy to an $n$-island system with migration
governed by graph~$G$. Here $T^n$ is the $n$-island monad: each
island carries its own PRNG state $S_i$, logs to its own
$\Lambda_i$, and shares a common configuration~$\Gamma$. Concretely,
$T^n\!A = \Gamma \to (S_1 \times \cdots \times S_n) \to
(A^n \times (S_1 \times \cdots \times S_n) \times
(\Lambda_1 \times \cdots \times \Lambda_n))$.
On objects,
$\mathcal{I}_G(\mathsf{Pop}(\Sigma)) = \mathsf{Pop}(\Sigma)^n$. On
morphisms, $\mathcal{I}_G(\sigma)(P_1, \ldots, P_n)$ applies
$\sigma$ to each island independently, then redistributes
individuals along edges of~$G$. Whether $\mathcal{I}_G$ preserves
composition depends on~$G$: when $G$ has no edges, composition is
preserved exactly (strict functor); when $G$ has edges, migration
after each composed step differs from migration after their
composition---the discrepancy is precisely the laxator.

\begin{definition}[Laxator]
\label{def:laxator}
Let $\mathcal{I}_G$ be the island functor for migration graph~$G$.
The \emph{laxator} of~$\mathcal{I}_G$ is the natural transformation
\[
  \varphi_G \;:\; \mathcal{I}_G(\sigma_1) \circ_K \mathcal{I}_G(\sigma_2)
  \;\Rightarrow\; \mathcal{I}_G(\sigma_1 \circ_K \sigma_2),
\]
measuring the discrepancy between composing island-level strategies
and lifting a composed strategy. When $G$ has no edges,
$\varphi_G = \mathrm{id}$ and $\mathcal{I}_G$ is a strict monoidal
functor. Whenever $G$ is connected, $\varphi_G \neq \mathrm{id}$,
making $\mathcal{I}_G$ lax.
\end{definition}

\begin{remark}[Status of the laxator]
\label{rem:laxator-status}
The explicit construction of $\varphi_G$---giving its components and
verifying coherence conditions---is the central open problem of this
framework. In this paper, we take a programmatic approach: we identify
the laxator as the key quantity, predict its behavior via spectral
theory (Section~\ref{sec:spectral}), and confirm the prediction
empirically (Section~\ref{sec:experiments}). The construction itself
is future work.
\end{remark}

%% ============================================================
%% Section 3: The Spectral Bridge
%% ============================================================
\section{The Spectral Bridge}
\label{sec:spectral}

The categorical framework identifies the laxator $\varphi_G$ as the
quantity governing diversity dynamics. We now connect its magnitude
to a computable graph invariant: the algebraic connectivity
$\lambda_2(G)$.

\subsection{Algebraic Connectivity and the Diversity Ordering}

\begin{conjecture}[Spectral prediction of topology ordering]
\label{prop:spectral}
Let $G$ be a connected migration graph on $n$ vertices.
The laxator magnitude $\|\varphi_G\|$ grows with
$\lambda_2(G)$. Since larger $\lambda_2$ implies faster mixing
and greater diversity erosion, the diversity ordering for
connected topologies is determined by $\lambda_2$: smaller
$\lambda_2$ yields higher sustained diversity.
\end{conjecture}

\paragraph{Heuristic argument.}
The spectral gap $\gamma$ of the lazy random walk on $G$ satisfies
$\gamma \geq \lambda_2(L)/(2\Delta)$ for maximum degree
$\Delta$~\cite{levin2009markov}, and $\gamma$ is monotone in
$\lambda_2$. The mixing time satisfies
$t_{\mathrm{mix}} = \Theta(\log(n/\varepsilon)/\gamma)$: smaller
$\lambda_2$ means slower mixing and higher sustained diversity.
Since the laxator measures the discrepancy introduced by
migration, and migration's homogenizing effect scales with
mixing rate, we expect $\|\varphi_G\|$ to grow with $\lambda_2$.
A formal proof requires the explicit construction of $\varphi_G$
(Remark~\ref{rem:laxator-status}); for now, this conjecture is
supported by the heuristic above and confirmed empirically in
Section~\ref{sec:experiments}.

The spectral values for our four connected topologies are:
\begin{itemize}
\item Ring $C_n$: $\lambda_2 = 2(1 - \cos(2\pi/n)) \approx 4\pi^2/n^2$
\item Star $K_{1,n-1}$: $\lambda_2 = 1$
\item Fully connected $K_n$: $\lambda_2 = n$
\end{itemize}
For random graphs, $\lambda_2$ depends on the realization.
Among connected topologies, the predicted ordering from smallest to
largest $\lambda_2$ (hence highest to lowest diversity) is:
ring~$>$~star~$>$~random~$>$~FC for $n \geq 7$.

\subsection{The \texorpdfstring{$n{=}5$}{n=5} Boundary}

At $n{=}5$, $\lambda_2(C_5) = 2(1 - \cos(2\pi/5)) = 1.382 > 1.0 =
\lambda_2(K_{1,4})$. Taken at face value, $\lambda_2$ alone
predicts star should preserve slightly more diversity than ring.
Empirically, ring edges out star in all six domains, but the
difference is not significant ($p = 0.14$, Section~\ref{sec:spectral-confirm}).
We interpret this as a \emph{scope condition}: when $\lambda_2$
values are close (here differing by only~$0.38$), second-order
effects---notably degree heterogeneity (the star hub sees $n{-}1$
edges while ring vertices see~$2$)---may dominate.
The robust prediction is that ring and star are
\emph{hard to distinguish} at $n{=}5$; the $\lambda_2$ ordering
becomes reliable only when the spectral gap between topologies
is large.

\subsection{The \texorpdfstring{$n{=}7$}{n=7} Prediction}

At $n{=}7$, $\lambda_2(C_7) = 2(1 - \cos(2\pi/7)) = 0.753 < 1.0 =
\lambda_2(K_{1,6})$. The inequality reverses: ring now couples
\emph{less strongly} than star. The theory makes a falsifiable
prediction: at $n \geq 7$, ring preserves strictly more diversity
than star. This is confirmed in Section~\ref{sec:spectral-confirm}.

\subsection{Categorical Interpretation}

The laxator magnitude grows with $\lambda_2$: smaller $\lambda_2$
means smaller laxator, closer to strict composition, more diversity
preserved. Among connected $k$-regular graphs, a \emph{Ramanujan}
graph maximizes $\lambda_2$ (Alon--Boppana bound)---maximally lax.
The diversity-optimal (minimally lax) connected topology minimizes
$\lambda_2$; for $k{=}2$ this is $C_n$, the ring.

The strict/lax dichotomy is thus a continuous family indexed by
$\lambda_2$, confirmed independently in oscillator
synchronization~\cite{sanz2026algebraic}.

%% ============================================================
%% Section 4: Experiments
%% ============================================================
\section{Experiments}
\label{sec:experiments}

\subsection{Setup}

All experiments: 5~islands $\times$ 16~individuals, 100~generations,
migration rate~$0.1$ every 5~generations, 30~seeds per
(topology, domain) pair. We sweep five topologies: none (no migration),
ring, star, random (Erd\H{o}s--R\'{e}nyi, resampled each migration
event), and fully connected. Diversity is measured as normalized
pairwise Hamming distance (binary genomes) or Euclidean distance
(real-valued genomes), consistent with Leinster's axiomatic
framework~\cite{leinster2021entropy}.

\subsection{Six Domains}

We select six domains spanning three axes: landscape structure
(monotone vs.\ intransitive), representation (fixed-length vs.\
variable-length), and constraint type (unconstrained, hard, soft).
The categorical structure---monad stack, pipeline shape, composition
operator---is identical across all domains; only genome type and
fitness function change.

\textbf{OneMax}~\cite{mitchell1992royal}: 20-bit bitstrings, fitness = number
of ones. The simplest possible landscape.
\textbf{Maze generation}: 60-bit binary vectors encoding wall
presence on a $6 \times 6$ grid; fitness combines BFS-derived path
length, dead-end ratio, and branching factor.
\textbf{Graph coloring}: 3-coloring of random graphs; fitness penalizes
constraint violations.
\textbf{Knapsack}: classic NP-hard optimization with capacity
constraints.
\textbf{No~Thanks!}: co-evolutionary card game; fitness is relative
to opponents within each island.
\textbf{Checkers}: 64-bit binary genomes encoding board-evaluation
weights; fitness determined by round-robin tournament within each
island, creating an intransitive landscape where fitness depends on
co-evolving opponents~\cite{samuel1959checkers}.

\subsection{Main Result: Universal Topology Ordering}

\begin{figure}[t]
\centering
\includegraphics[width=0.92\textwidth]{experiments/plots/multi_domain_topology_ordering.pdf}
\caption{Diversity (mean $\pm$ SE, 30~seeds, generation~99) across all six
  experimental domains: OneMax, Maze, Graph Coloring, Knapsack, No~Thanks!,
  and Checkers. The canonical ordering
  none~$>$~ring~$>$~star~$>$~random~$>$~FC is identical in every domain
  (Kendall's $W = 1.0$, $p = 0.00008$).
  Topology explains $28.7\times$ more variance than domain (two-way ANOVA).}
\label{fig:multi-domain-topology}
\end{figure}

\begin{observation}[Domain-independent topology ordering]
\label{obs:main}
Across six unrelated domains, the diversity ordering at
generation~99 is
none~$>$~ring~$>$~star~$>$~random~$>$~FC
with perfect concordance: Kendall's $W = 1.0$, $\chi^2 = 24.0$,
$p = 0.00008$; all 15 pairwise Spearman $\rho = 1.0$
(Figure~\ref{fig:multi-domain-topology}).
A two-way ANOVA confirms topology explains $28.7\times$ more
variance than domain.
\end{observation}

The none-to-ring transition produces a 35\% diversity drop---the
symmetry-breaking first coupling dominates. Each subsequent step
contributes at most 9\%. The laxator norm is continuous but
non-uniform: the first coupling is qualitatively different from
subsequent refinements.

Sorting network was tested as a seventh domain but is degenerate:
all topologies converge identically, identifying a scope condition
(when the fitness landscape imposes uniform convergence pressure,
topology effects vanish).

\subsection{Checkers: Co-evolutionary Stress Test}

Checkers is the most demanding test of our invariance claim.
Fitness is not a property of a genome but of a genome \emph{relative
to its co-evolving opponents}---strategy A may beat B, B may beat C,
yet C may beat A. Migration introduces not just genetic material but
new competitive context. Despite this, the ordering holds:
checkers shows the smallest phase transition ($11.1\%$
none-to-ring drop, consistent with co-evolutionary buffering where
migrant fitness depends on local context) yet preserves the full
five-topology ordering (ring vs.\ star: Cohen's $d = 0.577$,
$p = 0.029$; none vs.\ FC: $d = 5.76$).

\subsection{Spectral Predictions Confirmed}
\label{sec:spectral-confirm}

\paragraph{The $n{=}5$ boundary.}
At $n{=}5$, Conjecture~\ref{prop:spectral} predicts ring and star
should be hard to distinguish ($\lambda_2(C_5) = 1.382$ vs.\
$\lambda_2(K_{1,4}) = 1.0$, a gap of only~$0.38$). Confirmed:
Fisher's combined $p = 0.14$ across six domains---ring and star
diversities are statistically indistinguishable.
Note that ring edges out star consistently despite having
higher~$\lambda_2$; we attribute this to degree-heterogeneity
effects that dominate when $\lambda_2$ values are close
(Section~\ref{sec:spectral}).

\paragraph{The $n{=}7$ confirmation.}
At $n{=}7$, the spectral bridge predicts ring $>$ star in diversity
($\lambda_2(C_7) = 0.753 < 1.0 = \lambda_2(K_{1,6})$). A
confirmatory maze experiment (30~seeds, identical parameters) is
unambiguous: ring diversity $0.387 \pm 0.028$ versus star
$0.336 \pm 0.053$ ($p = 6.6 \times 10^{-5}$). Pure graph theory,
oscillator physics~\cite{sanz2026algebraic}, evolutionary graph
theory~\cite{brewster2025diversity}, and our experiments
converge---not just on the ordering, but on the \emph{boundary}
where it changes.

\paragraph{Time-varying topology and the $5.5\times$ inflation.}
For fixed graphs, snapshot $\lambda_2(G)$ perfectly predicts the
diversity ordering. For the time-varying random topology (resampled
each migration event), snapshot $\lambda_2$ fails:
$\bar{\lambda}_2 = 0.36$ is the lowest of all connected topologies,
yet random ranks fourth. The time-averaged adjacency matrix yields
effective $\lambda_2 = 39.74$ ($110\times$ the snapshot mean),
recovering the ordering. Categorically, fixed topology makes
$\mathcal{I}_G$ a strict endofunctor; time-varying topology produces
a genuinely lax endofunctor whose $5.5\times$ per-event diversity
inflation is the laxator's numerical signature.

\subsection{Fingerprints: Composition Determines Trajectory Shape}

\begin{definition}[Diversity fingerprint]
\label{def:fingerprint}
Let $\sigma : \mathsf{Pop}(\Sigma) \to T(\mathsf{Pop}(\Sigma))$ be a
strategy (a composed Kleisli morphism). The \emph{diversity fingerprint}
of~$\sigma$ is the sequence
$\mathbf{d}(\sigma) = (d_0, d_1, \ldots, d_n) \in \mathbb{R}^{n+1}$,
where $d_i$ is the genotypic diversity at generation~$i$.
\end{definition}

Beyond the topology ordering, we observe that strategies with
identical composition patterns produce matching diversity trajectory
\emph{shapes} across all tested domains: flat (monotonic decline),
hourglass (spike--crash--rebound), and island (stable oscillation
maintained by migration).
The categorical explanation: fingerprints are functions of the
composition pattern, which is invariant across domains. A formal
definition of ``shape equivalence'' remains open, but the empirical
pattern is consistent across all six domains.

%% ============================================================
%% Section 5: Discussion
%% ============================================================
\section{Discussion}
\label{sec:discussion}

\paragraph{Naturality interpretation.}
Why is the ordering domain-independent? Because it depends on
\emph{compositional structure}---preserved by the domain-change
functor---not on \emph{content} (genome type, fitness function).
$W = 1.0$ across six domains is \emph{consistent with} the topology
ordering being a natural transformation: when we replace checkers
genomes with maze genomes, the ordering is unchanged, just as a
natural transformation commutes with the underlying functor.
Making this precise requires formalizing the domain-change functor
and showing the ordering commutes with it---a goal for future work.

\paragraph{Practical implications.}
For EC practitioners: topology choice matters more than
domain-specific tuning. Ring is optimal among connected topologies.
At $n \geq 7$ islands, ring clearly outperforms star; at $n{=}5$,
they are interchangeable. The $28.7\times$ variance ratio suggests
that investing effort in topology selection yields greater returns
than investing in domain-specific diversity-preservation mechanisms.

\paragraph{Connections to other fields.}
Four independent lines confirm the $\lambda_2$ universality:
$\lambda_2$ governs synchronization onset in coupled
oscillators~\cite{sanz2026algebraic}, consensus time scales as
$N^3$ on cycles versus $N^2$ on complete
graphs~\cite{brewster2025diversity}, threshold graphs are globally
synchronizing~\cite{wu2025threshold}, and local structure determines
topologically protected
synchronization~\cite{ostilli2025topological}. Li et
al.~\cite{li2023emergence} formalize emergence via derived functors,
paralleling the laxator from a homological algebra perspective.
Ghani et al.~\cite{ghani2025approximate} show that exact equilibria
compose cleanly but approximate ones do not---the game-theoretic
analogue of our strict/lax dichotomy.

\paragraph{The general conjecture.}
We conjecture this pattern extends beyond EC: across optimization
paradigms, strict composition preserves compositional invariants while
lax composition degrades
them~\cite{chen2025scaling,kumar2026constitutional,
alphaevolve2026marl,alpay2025semantic,assumpcao2025codeevolve,
sanz2026algebraic,han2026conformity}.

\paragraph{Per-island asymmetry.}
At intermediate topologies (star, random), graph heterogeneity
produces per-island synchronization
asymmetry~\cite{zhou2026defect}; full time series are in the
companion repository.

\paragraph{Limitations.}
(a)~The laxator is identified but not explicitly constructed
(Remark~\ref{rem:laxator-status})---the central open problem.
(b)~Experiments are at $n{=}5$ and $n{=}7$; larger $n$ is
untested.
(c)~Sorting network is degenerate, identifying a scope condition:
domains with uniform convergence pressure may suppress topology
effects.
(d)~All experiments use homogeneous islands (same strategy on each
island).
(e)~At $n{=}5$, $\lambda_2$ predicts star~$>$~ring, but ring
edges out star consistently; degree heterogeneity likely introduces
a correction term not captured by $\lambda_2$ alone.

\paragraph{Future work.}
Three directions: (a)~explicit construction of the laxator
$\varphi_G$ and verification of coherence conditions;
(b)~heterogeneous islands with braided monoidal migration; and
(c)~scaling to larger $n$ to map the full $\lambda_2$ landscape.

%% ============================================================
%% Section 6: Conclusion
%% ============================================================
\section{Conclusion}
\label{sec:conclusion}

Migration topology determines diversity dynamics independently of
fitness landscape. The ordering
none~$>$~ring~$>$~star~$>$~random~$>$~fully connected holds with
perfect rank correlation across six fundamentally different
domains (Kendall's $W = 1.0$, $p = 0.00008$), with topology
explaining $28.7\times$ more variance than domain.

The mechanism: a categorical framework identifies GA operators as
Kleisli morphisms over a population monad, with the island functor's
laxator measuring migration-induced coupling. A spectral bridge
connects the laxator to $\lambda_2(G)$, predicting the diversity
ordering. The prediction is confirmed at $n{=}5$ (ring/star
indistinguishable, $p = 0.14$) and $n{=}7$ (ring clearly
outperforms star, $p = 6.6 \times 10^{-5}$)---the theory predicts
not just the ordering but the boundary where it changes.

Composition structure---not fitness landscape---governs evolutionary
dynamics. The explicit construction of the laxator, heterogeneous
islands, and formal proof of the general conjecture are the next
steps toward a complete categorical theory of evolutionary
computation.

%% ============================================================
%% Bibliography
%% ============================================================
{\small
\bibliographystyle{eptcs}
\bibliography{references}
}

\end{document}
